\subsection{THE HAUSDORFF SERIES}

Let \(\{U, V\}\) be a set with two elements.

DEFINITION 1. The element \(\mathbf{H} = \mathbf{U} \uplus \mathbf{V} = \log (\exp \mathbf{U} \cdot \exp \mathbf{V})\) (no. 2) of the Lie algebra \(\hat{\mathbf{L}}_{\mathbf{Q}}(\{\mathbf{U}, \mathbf{V}\})\) is called the Hausdorff series in the indeterminates \(\mathbf{U}\) and \(\mathbf{V}\) .

We denote by \(H_{n}\) (resp. \(H_{r,s}\) ) the homogeneous component of H of total degree n (resp. multidegree \((r, s)\) ). Then

\[
\mathrm{H} = \sum_ {n \geqslant 0} \mathrm{H} _ {n} = \sum_ {r, s \geqslant 0} \mathrm{H} _ {r, s}, \quad \mathrm{H} _ {n} = \sum_ {\substack {r + s = n \\ r, s \geqslant 0}} \mathrm{H} _ {r, s}.\tag{8}
\]

THEOREM 2. If \(r\) and \(s\) are two positive integers such that \(r + s \geqslant 1\) , then \(\mathrm{H}_{r,s} = \mathrm{H}_{r,s}' + \mathrm{H}_{r,s}''\) , where

\[
(r + s) \mathrm{H} _ {r, s} ^ {\prime} = \tag {9}
\]

\[
\sum_{m\geqslant 1}\frac{(-1)^{m - 1}}{m}\sum_{\substack{r_{1} + \dots +r_{m} = r\\ s_{1} + \dots +s_{m - 1} = s - 1\\ r_{1} + s_{1}\geqslant 1,\ldots ,r_{m - 1} + s_{m - 1}\geqslant 1}}\left(\left(\prod_{i = 1}^{m - 1}\frac{(\mathrm{adU})^{r_{i}}}{r_{i}!}\frac{(\mathrm{adV})^{s_{i}}}{s_{i}!}\right)\frac{(\mathrm{adU})^{r_{m}}}{r_{m}!}\right)(\mathrm{V})
\]

\[
(r + s) \mathrm{H} _ {r, s} ^ {\prime \prime} = \tag {10}
\]

\[
\sum_{m\geqslant 1}\frac{(-1)^{m - 1}}{m}\sum_{\substack{r_{1} + \dots +r_{m - 1} = r - 1\\ s_{1} + \dots +s_{m - 1} = s\\ r_{1} + s_{1}\geqslant 1,\dots ,r_{m - 1} + s_{m - 1}\geqslant 1}}\Bigg(\prod_{i = 1}^{m - 1}\frac{(\operatorname{adU})^{r_{i}}}{r_{i}!}\frac{(\operatorname{adV})^{s_{i}}}{s_{i}!}\Bigg)(U).
\]

In \(\hat{\mathbf{A}}_{\mathbf{q}}(\{\mathrm{U},\mathrm{V}\})\) , exp U .exp V = 1 + W, where W = \(\sum_{r+s\geqslant 1}\frac{\mathrm{U}^r}{r!}\frac{\mathrm{V}^s}{s!}\) whence H = \(\sum_{m\geqslant 1}(-1)^{m - 1}\mathrm{W}^m /m\) (no. 2), that is:

\[
\mathrm{H}_{r,s} = \sum_{m\geqslant 1}\frac{(-1)^{m - 1}}{m}\sum_{\substack{r_{1} + \dots +r_{m} = r\\ s_{1} + \dots +s_{m} = s\\ r_{1} + s_{1}\geqslant 1,\ldots ,r_{m} + s_{m}\geqslant 1}}\prod_{i = 1}^{m}\frac{\mathrm{U}^{r_{i}}}{r_{i}!}\frac{\mathrm{V}^{s_{i}}}{s_{i}!}.\tag{11}
\]

The linear mapping \(\mathbf{P}_n\) , defined by \(\mathbf{P}_n(x_1, \ldots, x_n) = \frac{1}{n} \left( \prod_{i=1}^{n-1} (\text{ad } x_i) \right)(x_n)\) for \(n \geqslant 1\) and \(x_1, \ldots, x_n\) in \(\{\mathrm{U}, \mathrm{V}\}\) , is a projector of \(\mathrm{A}_{\mathbf{Q}}^n(\{\mathrm{U}, \mathrm{V}\})\) onto \(\mathrm{L}_{\mathbf{Q}}^n(\{\mathrm{U}, \mathrm{V}\})\) (§ 3, no. 2, Corollary to Proposition 1); as \(\mathrm{H}_{r,s}\) belongs to \(\mathbf{L}_{\mathbf{Q}}^{r+s}(\{\mathrm{U}, \mathrm{V}\})\) , \(\mathrm{H}_{r,s} = \mathbf{P}_{r+s}(\mathrm{H}_{r,s})\) . Now

\[
\begin{array}{l} \mathrm{P} _ {r + s} \left(\prod_ {i = 1} ^ {m} \frac {\mathrm{U} ^ {r _ {i}} \mathrm{V} ^ {s _ {i}}}{r _ {i} ! s _ {i} !}\right) \\ = \frac {1}{r + s} \left(\left(\prod_ {i = 1} ^ {m - 1} \frac {(\operatorname{ad} \mathrm{U}) ^ {r _ {i}}}{r _ {i} !} \frac {(\operatorname{ad} \mathrm{V}) ^ {s _ {i}}}{s _ {i} !}\right) \frac {(\operatorname{ad} \mathrm{U}) ^ {r _ {m}}}{r _ {m} !} \frac {(\operatorname{ad} \mathrm{V}) ^ {s _ {m} - 1}}{s _ {m} !}\right) (\mathrm{V}) \end{array}\tag{12}
\]

when \(s_m \geqslant 1\) and

\[
\mathrm{P} _ {r + s} \left(\prod_ {i = 1} ^ {m} \frac {\mathrm{U} ^ {r _ {i}}}{r _ {i} !} \frac {\mathrm{V} ^ {s _ {i}}}{s _ {i} !}\right) = \frac {1}{r + s} \left(\left(\prod_ {i = 1} ^ {m - 1} \frac {(\mathrm{adU}) ^ {r _ {i}}}{r _ {i} !} \frac {(\mathrm{adV}) ^ {s _ {i}}}{s _ {i} !}\right) \frac {(\mathrm{adU}) ^ {r _ {m} - 1}}{r _ {m} !}\right) (\mathrm{U})\tag{13}
\]

when \(r_m \geqslant 1\) and \(s_m = 0\) . Moreover, obviously (ad \(t\) ) \(^{p-1}\) . \(t = 0\) if \(p \geqslant 2\) and (ad \(t\) ) \(^0\) . \(t = t\) . It follows that the two sides of (12) are zero when \(s_m \geqslant 2\) and those of (13) are zero when \(r_m \geqslant 2\) . The theorem then follows since \(H_{r,s}'\) is the sum of the terms of type (12) and \(H_{r,s}''\) is the sum of the terms of type (13).

Remarks. (1) We have defined (§ 3, no. 2, Remark) a projector Q of A(X) onto L(X) such that \(Q(a^{m}) = 0\) for \(a \in L(X)\) and \(m \geqslant 2\) and Q(1) = 0. Then H = Q(exp H) = Q(exp U.exp V), whence immediately

\[
\mathrm{H} _ {r, s} = \mathrm{Q} \left(\frac {\mathrm{U} ^ {r}}{r !} \frac {\mathrm{V} ^ {s}}{s !}\right) \quad \text { for } r + s \geqslant 1.\tag{14}
\]

\begin{enumerate}
\def\labelenumi{(\arabic{enumi})}
\setcounter{enumi}{1}
\tightlist
\item
  We have
\end{enumerate}

\[
\begin{array}{r l} \mathrm{H(U,V)} & \equiv \mathrm{U+V} + \frac {1}{2} [ \mathrm{U,V} ] + \frac {1}{12} [ \mathrm{U,[U,V]} ] \\ & \quad + \frac {1}{12} [ \mathrm{V,[V,U]} ] - \frac {1}{24} [ \mathrm{U,[V,[U,V]} ] ] \end{array} \tag {15}
\]

modulo \(\sum_{n\geqslant 5}\mathrm{L}^n (\{\mathrm{U},\mathrm{V}\})\)

\begin{enumerate}
\def\labelenumi{(\arabic{enumi})}
\setcounter{enumi}{2}
\tightlist
\item
  \(\mathrm{H}_{0,n} = \mathrm{H}_{n,0} = 0\) for every integer \(n\neq 1\) , whence
\end{enumerate}

\[
\mathrm{H} (\mathrm{U}, 0) = \mathrm{H} (0, \mathrm{U}) = \mathrm{U}.\tag{16}
\]

On the other hand, as \([U, -U] = 0\) ,

\[
\mathrm{H} (\mathrm{U}, - \mathrm{U}) = 0.\tag{17}
\]

